% GENERATED FILE. Edit canonical lessons, then run npm run generate:books.
% canonical-source-hash: fnv1a64:7df232e195650787
% canonical-lessons: JA-C284-senkyo, JA-C284-kokumin, JA-C284-shushou, JA-C284-arasou, JA-C284-warui

\chapter{Election Day}
\label{ch:ja-election-day}
\hlchaptermodality{284}

\begin{chapteropening}
\textbf{By the end of this chapter:} \emph{I can say who chose whom and how it went.}

Complete the last lesson of chapter 284: I can say who chose whom and how it went.
\end{chapteropening}

\section[\texorpdfstring{\ja{せんきょ} (senkyo)}{せんきょ (senkyo)}]{\ja{せんきょ} (senkyo) — an election}

\label{lesson:JA-C284-senkyo}



\begin{quote}
Before the new one: say the Japanese for information, then the Japanese for a fact.
\end{quote}

\subsection*{You'll want to know: \ja{せんきょ}}
\textbf{\ja{せんきょ}} — \emph{senkyo} — \textquotedblleft{}an election\textquotedblright{}.

\begin{quote}
\textbf{\ja{らいげつ} \ja{せんきょが} \ja{あります。}} — \emph{raigetsu senkyo ga arimasu} — \textquotedblleft{}There is an election next month.\textquotedblright{}
\end{quote}

\begin{cousinweb}
\textbf{\ja{せんきょ}} (\emph{senkyo}) is Sino-Japanese (kango): two Chinese characters read with their on'yomi, readings borrowed from Chinese. They are \emph{sen}, \textquotedblleft{}elect\textquotedblright{} (Middle Chinese \emph{sjwenX}) and \emph{kyo}, \textquotedblleft{}raise\textquotedblright{} (\emph{kjoX}) (KANJIDIC2; Middle Chinese in Baxter's transcription).
\end{cousinweb}

\begin{grammarlens}[title={What you've built}]
One more word for telling what happened.
\end{grammarlens}

\subsection*{Your turn}
\begin{itemize}
\raggedright
  \item \emph{Say it:} \emph{senkyo}
  \item \emph{Say it:} \emph{senkyo}, once more
  \item \emph{Say it:} \emph{raigetsu senkyo ga arimasu} — There is an election next month.
  \item \emph{Recall:} say the Japanese for information, then the Japanese for a fact, then say \emph{senkyo} again
\end{itemize}

\begin{tcolorbox}[breakable,colback=teal!4,colframe=teal!35!black,title={Before you move on}]
What does \ja{せんきょ} mean? (\textquotedblleft{}An election\textquotedblright{}.) Say it once more.
\end{tcolorbox}
\section[\texorpdfstring{\ja{こくみん} (kokumin)}{こくみん (kokumin)}]{\ja{こくみん} (kokumin) — the people (of a country); citizens}

\label{lesson:JA-C284-kokumin}



\begin{quote}
Before the new one: say the Japanese for a fact, then the Japanese for an election.
\end{quote}

\subsection*{You'll want to know: \ja{こくみん}}
\textbf{\ja{こくみん}} — \emph{kokumin} — \textquotedblleft{}the people (of a country); citizens\textquotedblright{}.

\begin{quote}
\textbf{\ja{せんきょで} \ja{こくみんが} \ja{せいふを} \ja{えらびます。}} — \emph{senkyo de kokumin ga seifu o erabimasu} — \textquotedblleft{}In an election the people choose the government.\textquotedblright{}
\end{quote}

\begin{cousinweb}
\textbf{\ja{こくみん}} (\emph{kokumin}) is Sino-Japanese (kango): two Chinese characters read with their on'yomi, readings borrowed from Chinese. They are \emph{koku}, \textquotedblleft{}country\textquotedblright{} (Middle Chinese \emph{kwok}) and \emph{min}, \textquotedblleft{}people\textquotedblright{} (\emph{mjin}) (KANJIDIC2; Middle Chinese in Baxter's transcription).
\end{cousinweb}

\begin{grammarlens}[title={What you've built}]
One more word for telling what happened.
\end{grammarlens}

\subsection*{Your turn}
\begin{itemize}
\raggedright
  \item \emph{Say it:} \emph{kokumin}
  \item \emph{Say it:} \emph{kokumin}, once more
  \item \emph{Say it:} \emph{senkyo de kokumin ga seifu o erabimasu} — In an election the people choose the government.
  \item \emph{Recall:} say the Japanese for a fact, then the Japanese for an election, then say \emph{kokumin} again
\end{itemize}

\begin{tcolorbox}[breakable,colback=teal!4,colframe=teal!35!black,title={Before you move on}]
What does \ja{こくみん} mean? (\textquotedblleft{}The people (of a country); citizens\textquotedblright{}.) Say it once more.
\end{tcolorbox}
\section[\texorpdfstring{\ja{しゅしょう} (shushō)}{しゅしょう (shushō)}]{\ja{しゅしょう} (shushō) — the prime minister}

\label{lesson:JA-C284-shushou}



\begin{quote}
Before the new one: say the Japanese for an election, then the Japanese for the people (of a country).
\end{quote}

\subsection*{You'll want to know: \ja{しゅしょう}}
\textbf{\ja{しゅしょう}} — \emph{shushō} — \textquotedblleft{}the prime minister\textquotedblright{}.

\begin{quote}
\textbf{\ja{けさ} \ja{しゅしょうが} \ja{はっぴょうしました。}} — \emph{kesa shushō ga happyō shimashita} — \textquotedblleft{}The prime minister made an announcement this morning.\textquotedblright{}
\end{quote}

\begin{cousinweb}
\textbf{\ja{しゅしょう}} (\emph{shushō}) is Sino-Japanese (kango): two Chinese characters read with their on'yomi, readings borrowed from Chinese. They are \emph{shu}, \textquotedblleft{}neck\textquotedblright{} (Middle Chinese \emph{syuwX}) and \emph{shō}, \textquotedblleft{}minister of state\textquotedblright{} (\emph{sjangH}) (KANJIDIC2; Middle Chinese in Baxter's transcription).
\end{cousinweb}

\begin{grammarlens}[title={What you've built}]
One more word for telling what happened.
\end{grammarlens}

\subsection*{Your turn}
\begin{itemize}
\raggedright
  \item \emph{Say it:} \emph{shushō}
  \item \emph{Say it:} \emph{shushō}, once more
  \item \emph{Say it:} \emph{kesa shushō ga happyō shimashita} — The prime minister made an announcement this morning.
  \item \emph{Recall:} say the Japanese for an election, then the Japanese for the people (of a country), then say \emph{shushō} again
\end{itemize}

\begin{tcolorbox}[breakable,colback=teal!4,colframe=teal!35!black,title={Before you move on}]
What does \ja{しゅしょう} mean? (\textquotedblleft{}The prime minister\textquotedblright{}.) Say it once more.
\end{tcolorbox}
\section[\texorpdfstring{\ja{あらそう} (arasou)}{あらそう (arasou)}]{\ja{あらそう} (arasou) — to compete; to contend}

\label{lesson:JA-C284-arasou}



\begin{quote}
Before the new one: say the Japanese for the people (of a country), then the Japanese for the prime minister.
\end{quote}

\subsection*{You'll want to know: \ja{あらそう}}
\textbf{\ja{あらそう}} — \emph{arasou} — \textquotedblleft{}to compete; to contend\textquotedblright{}.

\begin{quote}
\textbf{\ja{せんきょで} \ja{ふたりが} \ja{あらそっています。}} — \emph{senkyo de futari ga arasotte imasu} — \textquotedblleft{}Two candidates are fighting the election.\textquotedblright{}
\end{quote}

Polite \textbf{\ja{あらそいます}}, \emph{arasoimasu}; a godan verb. Past (polite): \textbf{\ja{ふたりが} \ja{あらそいました。}}, \emph{futari ga arasoimashita}.

\begin{cousinweb}
\textbf{\ja{あらそう}} (\emph{arasou}) is a native Japanese word (wago), inherited rather than borrowed from Chinese; its kanji was attached later as a kun'yomi (JMdict).
\end{cousinweb}

\begin{grammarlens}[title={What you've built}]
One more word for telling what happened.
\end{grammarlens}

\subsection*{Your turn}
\begin{itemize}
\raggedright
  \item \emph{Say it:} \emph{arasou}
  \item \emph{Say it:} \emph{arasou}, once more
  \item \emph{Say it:} \emph{senkyo de futari ga arasotte imasu} — Two candidates are fighting the election.
  \item \emph{Recall:} say the Japanese for the people (of a country), then the Japanese for the prime minister, then say \emph{arasou} again
\end{itemize}

\begin{tcolorbox}[breakable,colback=teal!4,colframe=teal!35!black,title={Before you move on}]
What does \ja{あらそう} mean? (\textquotedblleft{}To compete; to contend\textquotedblright{}.) Say it once more.
\end{tcolorbox}
\section[\texorpdfstring{\ja{わるい} (warui)}{わるい (warui)}]{\ja{わるい} (warui) — bad}

\label{lesson:JA-C284-warui}



\begin{quote}
Before the new one: say the Japanese for the prime minister, then the Japanese for to compete.
\end{quote}

\subsection*{You'll want to know: \ja{わるい}}
\textbf{\ja{わるい}} — \emph{warui} — \textquotedblleft{}bad\textquotedblright{}.

\begin{quote}
\textbf{\ja{あしたは} \ja{てんきが} \ja{わるいでしょう。}} — \emph{ashita wa tenki ga warui deshō} — \textquotedblleft{}The weather will probably be poor tomorrow.\textquotedblright{}
\end{quote}

An i-adjective.

\begin{cousinweb}
\textbf{\ja{わるい}} (\emph{warui}) is a native Japanese word (wago), inherited rather than borrowed from Chinese; its kanji was attached later as a kun'yomi (JMdict).
\end{cousinweb}

\begin{grammarlens}[title={What you've built}]
One more word for telling what happened.
\end{grammarlens}

\subsection*{Your turn}
\begin{itemize}
\raggedright
  \item \emph{Say it:} \emph{warui}
  \item \emph{Say it:} \emph{warui}, once more
  \item \emph{Say it:} \emph{ashita wa tenki ga warui deshō} — The weather will probably be poor tomorrow.
  \item \emph{Recall:} say the Japanese for the prime minister, then the Japanese for to compete, then say \emph{warui} again
\end{itemize}

\begin{tcolorbox}[breakable,colback=teal!4,colframe=teal!35!black,title={Before you move on}]
What does \ja{わるい} mean? (\textquotedblleft{}Bad\textquotedblright{}.) Say it once more.
\end{tcolorbox}
